Background only: the standard support-function formalism for Minkowski sums gives additivity of support functions and of exposed faces in a fixed direction; this is the convex-geometric framework behind the face description already available from Subproblems \(1\)–\(3\). I will use only those already-established face facts below. ([tropp.caltech.edu](https://tropp.caltech.edu/notes/Tro18-Lectures-Convex-LN.pdf))

**Proof.**

Fix \(i\). Set
\[
x_i:=s_i+t_i,\qquad x_{i+1}:=s_{i+1}+t_{i+1},\qquad E:=[x_i,x_{i+1}],
\]
and write \(\operatorname{relint}(E)\) for the relative interior of \(E\).

Recall that \(\sigma(t)\in\{\pm1\}\) is the sign attached to a translate \(t\in T\), while for \(p\in\mathbb{R}^2\),
\[
\sigma(p)=\sum_{\substack{t\in T\\ p-t\in S}}\sigma(t).
\]

From Subproblems \(1\)–\(3\), we may use the following facts.

1. \(E\) is an edge of \(P+Q\), and its direction is \(a_i=s_{i+1}-s_i\).
2. Every point of \(\operatorname{relint}(E)\) cancels:
   \[
   \sigma(p)=0\qquad\text{for all }p\in \operatorname{relint}(E).
   \]
3. If \(p\in E\) and \(p=t+s\) with \(t\in T\) and \(s\in S\), then necessarily
   \[
   t\in [t_i,t_{i+1}]
   \quad\text{and}\quad
   s\in [s_i,s_{i+1}].
   \]
   Since \(S\) is in convex position, the edge \([s_i,s_{i+1}]\) contains no point of \(S\) other than its endpoints, so
   \[
   [s_i,s_{i+1}]\cap S=\{s_i,s_{i+1}\}.
   \]
   Hence any contribution to a point of \(E\) comes from a translate in
   \[
   F_i=T\cap [t_i,t_{i+1}],
   \]
   and only via \(s_i\) or \(s_{i+1}\).

Because the face \([t_i,t_{i+1}]\) has direction \(a_i\), there is some real \(M\ge 0\) such that
\[
t_{i+1}=t_i+M a_i.
\]
Now define
\[
\Lambda:=\{\lambda\in[0,M]: t_i+\lambda a_i\in F_i\}.
\]
Then \(0,M\in\Lambda\).

For \(\lambda\in\Lambda\), write
\[
t(\lambda):=t_i+\lambda a_i.
\]
Observe that
\[
s_i+t(\lambda)=x_i+\lambda a_i,\qquad
s_{i+1}+t(\lambda)=x_i+(\lambda+1)a_i.
\]
So whenever \(0\le \lambda<M\), the point \(s_{i+1}+t(\lambda)\) lies in \(\operatorname{relint}(E)\), and whenever \(0<\lambda\le M\), the point \(s_i+t(\lambda)\) lies in \(\operatorname{relint}(E)\).

### Step \(1\): successor and predecessor relations

Let \(\lambda\in\Lambda\) with \(\lambda<M\), and set
\[
t:=t(\lambda),\qquad p:=s_{i+1}+t.
\]
Then \(p\in \operatorname{relint}(E)\), so \(\sigma(p)=0\).

One contribution to \(\sigma(p)\) is \(\sigma(t)\), coming from
\[
p-t=s_{i+1}.
\]
By the face description above, any other contribution must also come from a translate in \(F_i\), and the only other possible vertex of \(S\) is \(s_i\). But
\[
p-s_i=t+a_i.
\]
Therefore the only possible second translate is \(t+a_i\). Since \(\sigma(p)=0\), that second translate must actually be present, and
\[
\sigma(t+a_i)=-\sigma(t).
\]
Equivalently,
\[
\lambda+1\in\Lambda.
\tag{1}
\]

Similarly, let \(\lambda\in\Lambda\) with \(\lambda>0\), and set
\[
t:=t(\lambda),\qquad p:=s_i+t.
\]
Again \(p\in \operatorname{relint}(E)\), hence \(\sigma(p)=0\). One contribution is \(\sigma(t)\), and the only possible second one comes from
\[
p-s_{i+1}=t-a_i.
\]
So \(t-a_i\in F_i\) and
\[
\sigma(t-a_i)=-\sigma(t).
\]
Equivalently,
\[
\lambda-1\in\Lambda.
\tag{2}
\]

Thus every nonterminal point of \(F_i\) has a successor one step \(a_i\) ahead, and every noninitial point has a predecessor one step \(a_i\) behind.

### Step \(2\): no fractional chain can occur

I claim that
\[
\Lambda\subset \mathbb{Z}.
\]

Assume not. Pick some \(\lambda\in\Lambda\setminus\mathbb{Z}\), and let
\[
r:=\lfloor \lambda\rfloor,\qquad \lambda_0:=\lambda-r.
\]
By repeatedly applying \((2)\), we get
\[
\lambda_0\in\Lambda,
\qquad 0<\lambda_0<1.
\]
Now let
\[
t:=t(\lambda_0)=t_i+\lambda_0 a_i,\qquad p:=s_i+t.
\]
Since \(0<\lambda_0<1\), the point \(p=x_i+\lambda_0 a_i\) lies in \(\operatorname{relint}(E)\), so \(\sigma(p)=0\).

But by the same two-contributor analysis as above, the only possible translates that could contribute to \(p\) are \(t\) and \(t-a_i\). The second one cannot occur, because
\[
t-a_i=t_i+(\lambda_0-1)a_i
\]
has parameter \(\lambda_0-1<0\), so it is not in \([t_i,t_{i+1}]\), hence not in \(F_i\). Therefore \(p\) has exactly one contribution, namely \(\sigma(t)\), and thus
\[
\sigma(p)=\sigma(t)\in\{\pm1\},
\]
contradicting \(\sigma(p)=0\).

So indeed
\[
\Lambda\subset\mathbb{Z}.
\]

Since \(M\in\Lambda\), it follows that
\[
M\in\mathbb{Z}_{\ge 0}.
\]
Set
\[
m_i:=M.
\]

### Step \(3\): the face is the full arithmetic progression

We already know \(0\in\Lambda\). Repeatedly applying \((1)\) gives
\[
1,2,\dots,m_i\in\Lambda.
\]
Since also \(\Lambda\subset[0,m_i]\cap\mathbb{Z}\), we conclude
\[
\Lambda=\{0,1,\dots,m_i\}.
\]
Therefore
\[
F_i=\{t_i+k a_i:0\le k\le m_i\},
\qquad
t_{i+1}=t_i+m_i a_i.
\]

### Step \(4\): alternating signs

From the successor relation,
\[
\sigma\bigl(t_i+(k+1)a_i\bigr)=-\sigma\bigl(t_i+k a_i\bigr)
\qquad (0\le k<m_i).
\]
By induction on \(k\),
\[
\sigma(t_i+k a_i)=(-1)^k\,\sigma(t_i)
\qquad (0\le k\le m_i).
\]

This is exactly the required conclusion:
\[
F_i=\{t_i+k a_i:0\le k\le m_i\},
\qquad
t_{i+1}=t_i+m_i a_i,
\]
and the signs alternate along the chain. ∎